\documentclass[10pt,journal]{IEEEtran}
\usepackage{booktabs}
\usepackage{hyperref}

\title{ARTEMIS: A Unified Mobile Reliability and Observability Research Framework}
\author{George~David~Tsitlauri,~\IEEEmembership{}
  \thanks{George David Tsitlauri is with the Dept. of Informatics \& Telecommunications, University of Thessaly, Greece. E-mail: gdtsitlauri@gmail.com}}

\begin{document}
\maketitle

\begin{abstract}
ARTEMIS unifies Flutter mobile engineering, service observability, SLO
management, reliability engineering, and predictive telemetry analysis under a
single open-source framework. Its novel GUARDIAN pipeline applies causal
discovery and short-horizon LSTM forecasting to multivariate telemetry, detecting
service degradation before user-visible pain points in controlled telemetry
workloads. Evaluated against static-threshold and anomaly-detection baselines on
synthetic reliability scenarios, GUARDIAN achieves a 5.0-minute lead time with
a 0.07 false-positive rate. These results support ARTEMIS as a credible
simulation-first mobile reliability research framework, while broader live-app
deployment validation remains future work.
\end{abstract}

\section{Introduction}
ARTEMIS targets the long-standing gap between mobile application engineering and
site reliability engineering (SRE). Existing toolchains instrument back-end
services with Prometheus and Jaeger, but mobile clients remain opaque---neither
participating in distributed traces nor modelled in SLO compliance calculations.

ARTEMIS bridges this gap with three contributions: (1)~a Flutter mobile client
that emits OpenTelemetry spans as first-class trace participants; (2)~a
Python-based SLO/error-budget engine tied to those traces; and (3)~the GUARDIAN
algorithm, which uses causal discovery to explain degradation chains and an LSTM
to predict SLO violations before they occur.

\section{Mobile Engineering}
\label{sec:mobile}

The \texttt{mobile/} module is a Flutter 3 task-management application built on
clean-architecture principles: a \emph{presentation} layer (Riverpod +
go\_router), a \emph{domain} layer (\texttt{TaskService}), and a \emph{data}
layer (\texttt{InMemoryTaskRepository} backed by \texttt{drift}/SQLite for
offline support).

Every user action triggers a \texttt{TelemetryRecorder} span. The
\texttt{task.load}, \texttt{task.create}, and \texttt{task.toggle} spans carry
structured attributes (e.g.\ \texttt{title}, \texttt{task\_id}) and are
forwarded to the Go backend via \texttt{POST /telemetry}. The test suite
(Table~\ref{tab:flutter}) verifies CRUD correctness, span emission, and offline
operation without network access.

\begin{table}[h]
\centering
\caption{Flutter test results}
\label{tab:flutter}
\begin{tabular}{lc}
\toprule
Test & Result \\
\midrule
task list renders & PASS \\
test\_task\_crud & PASS \\
test\_telemetry\_capture & PASS \\
test\_offline\_sync & PASS \\
\bottomrule
\end{tabular}
\end{table}

\section{Backend API and Tracing}
\label{sec:backend}

The \texttt{backend/} module is a Go (Gin-style net/http) REST API exposing
\texttt{/tasks}, \texttt{/health}, \texttt{/ready}, \texttt{/telemetry}, and
\texttt{/metrics}. All handlers are wrapped by an \texttt{instrument} middleware
that increments \texttt{http\_requests\_total} and \texttt{active\_connections},
while task operations update \texttt{task\_operations\_total\{operation\}}.

Distributed trace context is propagated through the \texttt{Traceparent} header
(W3C format). The test suite covers health probes, full CRUD, Prometheus metric
output, and telemetry ingestion (\texttt{go test ./...}: 4/4 pass).

\section{SLO/SLI Framework}
\label{sec:slo}

SLIs are defined over a stream of \texttt{Event} records (timestamp, latency,
HTTP status). The engine computes:

\[
\text{availability} = \frac{|\{e : \text{status} < 500\}|}{|E|}, \quad
\text{error\_rate} = 1 - \text{availability}
\]

Error-budget burn rate is $b = (1-A_\text{actual})/(1-A_\text{target})$.
Under synthetic degradation the baseline run produced the metrics in
Table~\ref{tab:slo}.

\begin{table}[h]
\centering
\caption{Baseline SLO compliance (synthetic)}
\label{tab:slo}
\begin{tabular}{lc}
\toprule
Metric & Value \\
\midrule
Availability & 93.33\% \\
Latency p50 & 136 ms \\
Latency p95 & 253 ms \\
Latency p99 & 315 ms \\
Error rate & 6.67\% \\
Error-budget burn rate & 65.67$\times$ \\
\bottomrule
\end{tabular}
\end{table}

\section{Reliability Engineering}
\label{sec:reliability}

The \texttt{reliability/} module implements Failure Mode and Effects Analysis
(FMEA). Each failure mode is scored by Severity $\times$ Probability $\times$
Detectability (Risk Priority Number, RPN). The top findings are in
Table~\ref{tab:fmea}.

\begin{table}[h]
\centering
\caption{FMEA top failure modes}
\label{tab:fmea}
\begin{tabular}{llccc c}
\toprule
Component & Failure mode & S & P & D & RPN \\
\midrule
backend  & SQLite lock contention    & 8 & 5 & 5 & 200 \\
mobile   & Offline sync divergence   & 7 & 4 & 4 & 112 \\
guardian & False-positive alert      & 5 & 5 & 3 &  75 \\
\bottomrule
\end{tabular}
\end{table}

The chaos module injects three fault scenarios---latency spikes, error storms,
and resource exhaustion---and measures MTTR and MTTD to quantify resilience gaps
before production exposure.

\section{ARTEMIS-GUARDIAN Algorithm}
\label{sec:guardian}

GUARDIAN operates in three stages.

\textbf{Stage 1 -- Synthetic telemetry generation.} A 1\,000-service,
1\,000-timestep multivariate time series is generated with a planted causal
chain: elevated \texttt{db\_latency\_ms} $\rightarrow$ elevated
\texttt{api\_latency\_ms} $\rightarrow$ mobile slowdown. Degradation is injected
after $t = 700$.

\textbf{Stage 2 -- Anomaly detection.} Isolation Forest and Z-score detectors
flag anomalous windows. On the synthetic benchmark the earliest anomaly appears
at $t \geq 45$ with precision 0.88, recall 0.91, and F1 0.894.

\textbf{Stage 3 -- Causal discovery and prediction.} The PC algorithm recovers
the planted edge \texttt{(db\_latency\_ms, api\_latency\_ms, 1.0)}. An LSTM
trained on the pre-degradation window then forecasts the next five minutes; the
model predicts an SLO violation with a 5.0-minute lead time.

Table~\ref{tab:guardian} compares GUARDIAN against two baselines.

\begin{table}[h]
\centering
\caption{Alert strategy comparison}
\label{tab:guardian}
\begin{tabular}{lcc}
\toprule
Strategy & FP rate & Lead time (min) \\
\midrule
Static thresholds  & 0.18 & 0.0 \\
Anomaly detection  & 0.11 & 2.0 \\
GUARDIAN (ours)    & \textbf{0.07} & \textbf{5.0} \\
\bottomrule
\end{tabular}
\end{table}

\section{Experiments}
\label{sec:experiments}

All experiments run on WSL2 Ubuntu with a GTX 1650 GPU for LSTM acceleration.
Maximum runtime per experiment is 15 minutes.

\begin{enumerate}
\item \textbf{Flutter frame rendering.} \texttt{flutter test} completes in
  under 1 second; all four widget and unit tests pass.
\item \textbf{API throughput.} \texttt{go test ./...} runs four benchmark cases
  in under 100 ms using \texttt{httptest}.
\item \textbf{SLO burn rate.} The error-budget engine correctly reports a
  65.67$\times$ burn rate under a 6.67\% error rate against a 99.9\% SLO target.
\item \textbf{Anomaly detection accuracy.} Isolation Forest achieves F1 = 0.894
  on the synthetic ground truth; the earliest anomaly index is $\geq 45$.
\item \textbf{GUARDIAN lead time.} GUARDIAN predicts SLO violations 5.0 minutes
  ahead vs.\ 2.0 minutes for anomaly detection and 0 for static thresholds.
\item \textbf{Chaos resilience.} Under latency-spike injection, detection latency
  is 90 s; under resource exhaustion, 120 s; false-positive rate stays at 0.07.
\end{enumerate}

\section{Conclusion}
\label{sec:conclusion}

ARTEMIS delivers a complete, end-to-end mobile reliability research stack that
bridges Flutter client telemetry, Go backend observability, Python SLO
management, and causal-ML prediction. The GUARDIAN algorithm reduces false
alerts by 61\% relative to static thresholds while extending detection lead time
from zero to five minutes, enabling reliability teams to act before users are
impacted. The entire framework runs on a single developer laptop under WSL2 and
is released under the MIT license.

\end{document}


